\begin{center}
    \section*{\fontsize{20}{20}\selectfont Chapter 2}
\end{center}
\vspace{10mm}
\section{Background Study and Literature Review}

Brain tumour segmentation has advanced substantially over the past two decades, driven by progress in both medical imaging technology and deep learning methodology. Accurately delineating tumour boundaries from multi-parametric MRI scans is fundamental to diagnosis, treatment planning, and longitudinal patient monitoring. This chapter surveys the evolution of automated segmentation methods — from classical image processing techniques through the convolutional neural network era, to emerging graph-based approaches — and identifies the persistent challenges of computational efficiency and clinical deployability that motivate the GNN framework proposed in this thesis.

\subsection{Evolution of Medical Image Segmentation}

\subsubsection{Classical Approaches}

Early brain tumour segmentation relied on traditional computer vision techniques including intensity thresholding, region growing, and atlas-based registration. These methods required extensive manual parameter tuning and struggled to accommodate the high inter-patient variability in tumour appearance. The BraTS Challenge, established in 2012 \cite{menze2015multimodal}, standardised evaluation protocols and provided large-scale annotated datasets, substantially accelerating innovation in automated segmentation.

\subsubsection{The CNN Revolution: U-Net and Its Variants}

Fully convolutional networks fundamentally transformed medical image segmentation. Ronneberger et al.\ \cite{ronneberger2015unet} introduced U-Net in 2015, establishing the symmetric encoder-decoder architecture with skip connections that preserves spatial information during upsampling. U-Net became the standard baseline for medical segmentation, substantially outperforming classical methods while requiring comparatively little training data through the use of aggressive data augmentation.

\c{C}i\c{c}ek et al.\ \cite{cicek20163d} extended U-Net to 3D medical imaging through volumetric convolutions, enabling direct processing of multi-slice MRI volumes. However, 3D U-Net's computational requirements scale cubically with volume dimensions, necessitating 8GB+ of VRAM and limiting deployment on resource-constrained hardware.

Isensee et al.\ \cite{isensee2021nnunet} introduced nnU-Net, a self-configuring framework that automatically optimises preprocessing, architecture, and training hyperparameters for diverse medical segmentation tasks, achieving state-of-the-art results across multiple datasets. However, automatic configuration requires extensive computational resources for hyperparameter search, and the resulting models remain large (80M+ parameters).

\subsubsection{Transformer-Based Architectures}

Recent works have explored vision transformers for medical segmentation. Hatamizadeh et al. \cite{hatamizadeh2022swin} proposed Swin-UNETR, combining Swin Transformer encoders with CNN decoders, achieving 93.8\% Dice on BraTS 2021. TransBTS \cite{wang2021transbts} similarly integrates transformer self-attention with 3D convolutions. While these hybrid architectures push accuracy boundaries, they require even larger computational budgets (92M+ parameters, 24GB+ VRAM), exacerbating deployment challenges for clinical settings.

\subsection{Graph Neural Networks in Medical Imaging}

\subsubsection{Fundamentals of Graph Representation Learning}

Graph neural networks extend deep learning to non-Euclidean data structures, learning representations through iterative message passing between connected nodes \cite{hamilton2017graphsage, kipf2017gcn}. Hamilton et al. \cite{hamilton2017graphsage} introduced GraphSAGE (Graph Sample and Aggregate), which learns node embeddings by sampling and aggregating features from local neighborhoods. The inductive learning capability of GraphSAGE enables generalization to unseen graph structures, making it suitable for medical imaging where each patient represents a distinct graph.

\subsubsection{Superpixel Segmentation}

Superpixel algorithms group perceptually similar pixels into compact regions, providing mid-level representations that preserve object boundaries while reducing dimensionality. Achanta et al. \cite{achanta2012slic} developed Simple Linear Iterative Clustering (SLIC), a computationally efficient method that performs k-means clustering in combined color-spatial space. SLIC generates uniform superpixels with low computational cost (linear in pixel count) and controllable compactness, making it ideal for large-scale medical image preprocessing.

\subsubsection{Graph-Based Medical Segmentation: Existing Work}

Several studies have explored graph representations for medical imaging:

\textbf{Hybrid CNN-GNN Approaches:} Most existing work combines CNNs for feature extraction with GNNs for spatial reasoning. These hybrid methods retain the computational overhead of volumetric convolutions while adding graph processing costs, yielding limited efficiency gains.

\textbf{Region Adjacency Graphs:} Some methods construct graphs where nodes represent anatomical structures (organs, lobes) and edges encode spatial relationships. While interpretable, these approaches require prior segmentation of regions, introducing preprocessing dependencies.

\textbf{Point Cloud Representations:} Voxel-to-point sampling combined with PointNet-style architectures has been investigated, but point clouds discard spatial structure information that graphs can preserve through explicit edge connectivity.

\subsection{Research Gap and Motivation}

\subsubsection{The Efficiency-Accuracy Dilemma}

Despite impressive accuracy, modern deep learning methods face persistent deployment barriers:

\textbf{Computational Requirements:} State-of-the-art models including nnU-Net, Swin-UNETR, and TransBTS require 68M--92M parameters and 8--24GB of GPU memory, limiting deployment to high-end research hardware. Rural clinics, mobile diagnostic units, and healthcare systems in developing regions typically lack access to such infrastructure.

\textbf{Inference Latency:} Processing a 240$\times$240$\times$155 MRI volume through deep 3D CNNs requires 50--100ms per patient, which is problematic for real-time clinical workflows where immediate feedback is expected.

\textbf{Energy Consumption:} Large models consume substantial power during inference. As medical AI scales to millions of patients globally, the environmental impact becomes significant, motivating the development of architectures that balance accuracy with energy efficiency.

\subsubsection{Unexplored GNN Potential}

Graph neural networks have demonstrated strong performance in social network analysis, molecular modelling, and recommendation systems. Their application to medical image segmentation, however, remains relatively unexplored:

\textbf{Pure GNN Approaches Absent:} Most medical imaging GNN research employs hybrid CNN-GNN designs in which a convolutional backbone provides features and a GNN provides spatial reasoning. No prior work has rigorously benchmarked a pure, lightweight GNN architecture — without a CNN backbone — for efficient brain tumour segmentation.

\textbf{Superpixel-GNN Integration Underexplored:} Combining superpixel-based preprocessing (enabling up to 890$\times$ dimensionality reduction) with inductive graph learning via GraphSAGE is an underexplored avenue that may achieve competitive accuracy with dramatic efficiency improvements.

\textbf{Reproducibility Gap:} Medical AI research has faced scrutiny over reproducibility — insufficient documentation, unreported hyperparameters, and potential data leakage have undermined confidence in reported results. Establishing rigorous validation standards including patient-level splits and comprehensive integrity checks is critical for the field's maturity.

\subsection{Positioning This Work}

This thesis addresses the identified research gap through three key contributions:

\textbf{1. Pure GNN Architecture:} We demonstrate that a lightweight GraphSAGE model (439K parameters) without any convolutional components can achieve 92.92\% Dice—approaching state-of-the-art transformers while offering 156$\times$ parameter reduction and 6.9$\times$ inference speedup.

\textbf{2. Rigorous Validation Framework:} Patient-level stratified cross-validation with 15 comprehensive integrity checks ensures zero data leakage and establishes reproducibility standards for the field.

\textbf{3. Clinical Deployment Focus:} By achieving real-time performance (12.7ms per patient) on consumer-grade hardware (RTX 2060 6GB), we demonstrate practical viability for resource-constrained healthcare settings.

\vspace{5mm}

This literature review establishes that while CNNs and transformers have pushed accuracy boundaries for brain tumour segmentation, their computational demands create deployment barriers that limit clinical impact. Graph neural networks, through exploitation of tumour sparsity via superpixel representations, offer a promising yet underexplored paradigm for efficient medical image analysis. The following chapters detail the systematic investigation of this approach, demonstrating that competitive accuracy does not require massive models when architectural design aligns with the structural properties of the target domain.

\clearpage